\documentclass{article}
\usepackage[italian]{babel}
\usepackage{float}
\usepackage{graphicx}
\usepackage[a4paper,top=2cm,bottom=2cm,left=3cm,right=3cm,marginparwidth=1.75cm]{geometry}
\usepackage{amsmath}
\usepackage{esdiff}
\usepackage{amsfonts}
\usepackage{amssymb}
\usepackage[colorlinks=true, allcolors=blue]{hyperref}
\setlength{\parskip}{\baselineskip}%
\setlength{\parindent}{0pt}%
\newcommand{\bigslant}[2]{{\raisebox{.1em}{$#1$}\left/\raisebox{-.2em}{$#2$}\right.}}
\newcommand{\norm}[1]{\left\lVert#1\right\rVert}

\title{Documentazione Progetto InfoLab}
\author{Giovanni Menicucci}
\date{June 2026}

\begin{document}

\maketitle

\section{Introduzione}
Questo documento tiene traccia dello sviluppo del progetto per il corso di Informatica con Laboratorio 2025/26, riguardante la soluzione dell'equazione del calore sul quadrato $[0,1]^2$.

\section{Struttura del progetto}
I task richiesti nel progetto vengono eseguiti con codici C++ (reperibili nella cartella "codice\_cpp" o Python (nella cartella "codice\_python").
Alcuni codici producono file di testo (nella cartella "dati"), che poi vengono letti dai codici successivi.
Tutti i task sono ripetuti per $N = 32, 64,128,256,512,1024$ (e anche per $N = 4$ che è verificabile a mano), e se necessario per i due ordinamenti definiti nel testo.  

Le sezioni seguenti contengono i moduli principali che compongono il progetto.

\section{Discretizzazione del dominio}
Input: numero naturale $N$.
Il dominio $[0,1]^2$ è discretizzato con una griglia equispaziata di $(N+2) \times (N+2)$ punti.

Output: grafo di adiacenza, i cui nodi sono i punti interni alla griglia, con archi tra nodi primi vicini nella griglia.

File coinvolti: codice $\texttt{discretizzazione.cpp}$ che genera un file $\texttt{coords.txt}$ che rappresenta la griglia e un file $\texttt{connectivity.txt}$ che rappresenta il grafo.

Strutture dati necessarie: 
\begin{itemize}
\item $\texttt{struct}$ per i nodi, per memorizzare indice progressivo $n$, indici della griglia $i,j$, coordinate spaziali $x,y$.
\item $\texttt{vector}$ per contenere la sequenza dei nodi interni, ordinati secondo l'indice $n$, per poi esportarli su $\texttt{coords.txt}$.
\item $\texttt{vector}$ di $\texttt{vector}$ per la matrice $\texttt{griglia}$ che associa alle coordinate spaziali di ciascun nodo interno il relativo indice progressivo, e associa -1 ai nodi di bordo.
\end{itemize}

Funzioni:
\begin{itemize}
\item $\texttt{GeneraGriglia}$.
Calcola il passo della griglia e genera i singoli nodi interni con un doppio ciclo for che assegna indice progressivo $n$ (ordine lessicografico), coordinate nella griglia e spaziali.
Riempie il vettore dei nodi interni e la matrice $\texttt{griglia}$, aumentando $n$ di 1 alla fine di ogni iterazione.
La matrice è inizializzata nel $\texttt{main}$ con valori $-1$.
Complessità: $O(N^2)$, a causa del doppio ciclo for.
All'interno di ogni iterazione si fanno solo operazioni elementari.

\item $\texttt{EsportaPuntiInterni}$.
Crea il file $\texttt{coordsN.txt}$, scorre sul vettore dei nodi interni e li scrive sul file, usando una riga per ciascun nodo, contenente $n,i,j,x,y$.
Complessità: $O(N^2)$ (lunghezza del vettore su cui scorro).

\item $\texttt{GeneraEsportaGrafo}$.
Crea il file $\texttt{connectivityN.txt}$.
Svolge un doppio ciclo for, sugli indici dei nodi interni.
Ad ogni iterazione, usa la matrice $\texttt{griglia}$ per verificare se i 2 vicini a destra e in alto del nodo corrente sono di bordo o no (se un nodo è di bordo, il valore assegnatogli nella matrice è rimasto -1, cioè non è stato sovrascritto con l'indice progressivo).
Controlla solo i vicini a destra e in alto, perché se controllassi tutti e 4 conterei due volte gli archi.
Se il vicino non è di bordo, scrive in una riga del file l'indice progressivo $e$ dell'arco (cioè quanti archi ho creato fino a questo) e gli indici progressivi dei due nodi.
Complessità: $O(N^2)$, a causa del doppio ciclo for.

\end{itemize}

Complessità totale: $O(N^2)$.


\section{Riordinamento dei nodi}
Input: nodi definiti sopra.

Output: nodi riordinati con nested dissection.

File coinvolti: codice $\texttt{riordinamento.cpp}$ che prende il file $\texttt{coords.txt}$ generato nel modulo precedente, e genera un file $\texttt{ordering.txt}$ che rappresenta la permutazione dei nodi.

Strutture dati necessarie:
\begin{itemize}
\item $\texttt{struct}$ per i nodi, come nel modulo di discretizzazione.
\item $\texttt{vector}$ per le sequenze di nodi.
\end{itemize}

Funzioni:
\begin{itemize}
\item $\texttt{GeneraVettoreNodi}$. Prende in input $N$ e un vettore vuoto definito nel main.
Scorre sul file $\texttt{coordsN.txt}$ e legge gli indici $n,i,j,x,y$ dei nodi, riempiendo il vettore con i nodi.
In pratica inverte quanto fatto nella funzione $\texttt{EsportaPuntiInterni}$ del modulo di discretizzazione.
Complessità: $O(N^2)$, che è il numero di nodi.

\item $\texttt{GeneraGriglia}$.
Svolge un doppio ciclo for che assegna indice progressivo $n$ agli elementi $(i,j)$ della matrice $\texttt{griglia}$.
La matrice è inizializzata nel $\texttt{main}$ con valori $-1$.
Complessità: $O(N^2)$, a causa del doppio ciclo for.

Questa funzione è analoga alla funzione omonima del modulo precdente. In questo modulo viene usata dalla funzione $\texttt{GeneraNuovoOrdinamentoVeloce}$.

\item $\texttt{GeneraNuovoOrdinamento}$.
Prende in input il vettore dei nodi riempito dalla funzione $\texttt{GeneraVettoreNodi}$, un vettore vuoto che conterrà i nodi riordinati, e gli indici $i,j$ massimo e minimo dei nodi contenuti nel rettangolo da riordinare (quando la chiamo dal main, sarà $i_{min} = j_{min} = 1$, $i_{max} = j_{max} = N$, cioè tutta la griglia dei nodi interni).

Calcola le dimensioni $N_1, N_2$ del rettangolo da dividere in tre.
Il caso base della ricorsione è quando il rettangolo non è divisibile ulteriormente, cioè contiene un solo nodo (oppure non è nemmeno definito perché $i_{min} >i_{max}$, può succedere).
Se si supera il caso base, si divide il rettangolo verticalmente o orizzontalmente, in base a quale lato del rettangolo è maggiore.
Si prende la riga o colonna di mezzo come separatore e si scorre sul vettore dei nodi, assegnandoli ai tre vettori $V_1, V_2, V_S$ definiti nel testo.
Poi applica la funzione a $V_1$ e $V_2$ (passo ricorsivo), facendo attenzione agli indici che li delimitano.
Infine aggiunge il vettore $V_S$ in fondo al vettore dei nodi riordinati; gli altri nodi (che non fossero nei casi base) sono stati aggiunti da questa stessa riga di codice durante la ricorsione.

Complessità: $O(N^2 \log N )$, perché ad ogni passo della ricorsione devo scorrere sul vettore dei nodi corrente, che ha lunghezza $O(N^2)$, e ad ogni passo della ricorsione divido per 2 (circa) il lato del rettangolo considerato, quindi il numero di passi è $O(\log N)$.

\item $\texttt{GeneraNuovoOrdinamentoVeloce}$.
Come la versione precedente prende in input il vettore dei nodi riempito dalla funzione $\texttt{GeneraVettoreNodi}$, un vettore vuoto che conterrà i nodi riordinati, e gli indici $i,j$ massimo e minimo dei nodi contenuti nel rettangolo da riordinare.
In più prende in input la griglia definita dalla funzione $\texttt{GeneraGriglia}$.

Svolge la stessa ricorsione della funzione precedente, ma evita di scorrere sul vettore dei nodi per decidere quali appartengono a $V_1$, $V_2$, $V_S$, dato che in effetti possiamo capire a priori come dividere i nodi, usando gli indici.
Perciò esegue subito le due chiamate ricorsive sugli insiemi $V_1$ e $V_2$, che adesso non vengono salvati ma sono identificati solo dagli indici che li delimitano.
Infine aggiunge $V_S$ (anch'esso non salvato) in fondo al vettore dei nodi riordinati.

Complessità: $O(N^2)$. All'interno di ogni chiamata della ricorsione, se il numero di nodi coinvolti è $n$, si scorre su $V_s$ che ha lunghezza $O(\sqrt{n})$. Si parte da $n_0 = N^2$ e ad ogni passo si dividono i nodi circa in due parti uguali a cui si applica di nuovo l'algoritmo, quindi il tempo necessario è dell'ordine di:
$$\sum_{k=0} ^{\log_2 n_0} 2^k \sqrt{\frac{n_0}{2^k}} = \sqrt{n_0} \sum_{k=0} ^{\log_2 n_0} 2^{k/2} = \sqrt{n_0} \frac{2^{(\log_2 n_0 +1)/2} - 1}{2^{1/2} - 1} = O(\sqrt{n_0} \cdot 2^{(\log_2 n_0 +1)/2}) = O(n_0).$$

\item $\texttt{EsportaOrdinamento}$.
Crea il file $\texttt{ordering.txt}$.
Scorre sul vettore dei nodi riordinati, generato da una delle due funzioni sopra.
Per ciascun nodo, scrive in una riga del file la posizione $m$ nel nuovo ordine e l'indice progressivo naturale $n$ contenuto nella struct.
Complessità: $O(N^2)$, che è la lunghezza del vettore.
\end{itemize}

Complessità totale: $O(N^2 \log N)$ se si usa il riordinamento lento, $O(N^2)$ con quello veloce.


\section{Approssimazione dell'equazione differenziale con sistema lineare}
Input: dominio dell'equazione differenziale discretizzato (ordinato nei due modi), coefficiente $\kappa$, funzione sorgente, condizione al bordo.

Output: matrice e termine noto del sistema lineare che approssima l'equazione.

File coinvolti: codice $\texttt{approx\_lineare.cpp}$ che prende i file $\texttt{coords.txt}$ (generato nel modulo "discretizzazione") e $\texttt{ordering.txt}$ (generato nel modulo "riordinamento") e  genera un file $\texttt{A.txt}$ che rappresenta la matrice del sistema lineare, e un file $\texttt{rhs.txt}$ che rappresenta il termine noto.

Strutture dati necessarie: 
\begin{itemize}
\item $\texttt{vector}$ per matrici e vettori.
\item $\texttt{pair}$ per coppie di indici.
\end{itemize}

Funzioni:
\begin{itemize}
\item $\texttt{GeneraIndiciBordo}$.
Per accettare in input una condizione al bordo non identicamente nulla, assegno a ogni punto del bordo della griglia un indice progressivo tra $-1$ e $-4N-4$, cioè creo una funzione biunivoca $p(i,j)$ tra gli indici $(i,j)$ di bordo e gli interi $\{-1, \dots, -4N-4\}$.
Lo faccio scorrendo in senso orario sul bordo, partendo dall'origine.
Assegno ai punti del bordo indici negativi per poter distinguere facilmente punti di bordo e non.
Scrivo gli indici assegnati ai punti del bordo nelle corrispondenti entrate di una matrice $(N+2) \times (N+2)$.

Se nella funzione $\texttt{GeneraEsportaMatriceRhs}$ (definita sotto) si vuole prendere in input una condizione al bordo, essa dovrà essere un vettore $v$ di $4N+4$ numeri reali, ordinati in modo che $u_0(x_i,y_j) = v[-p(i,j) - 1]$.

Complessità: $O(N)$, che è la lunghezza del bordo.

\item $\texttt{GeneraMappaNaturale}$.
Genera una mappa da $(i,j)$ a $n$ (ordine naturale), e la sua inversa.
Legge il file $\texttt{coords.txt}$ e scrive su una matrice $(N+2)\times (N+2)$ (che sarà la stessa usata per i punti di bordo) l'indice progressivo $n$ associato a ogni valore di $(i,j)$ interno.
Inoltre riempie un vettore che fa l'inverso, cioè associa la coppia $(i,j)$ a $n$ (solo per i punti interni).
Infine chiama la funzione $\texttt{GeneraIndiciBordo}$ per associare i valori da $-1$ a $-4N-4$ ai punti di bordo.

Complessità: $O(N^2)$, il numero di elementi nella matrice..

\item $\texttt{GeneraMappaOrdinata}$.
Genera una mappa da $(i,j)$ a $m$ (riordinamento), e la sua inversa.
Legge il file $\texttt{ordering.txt}$ e fa l'analogo della funzione precedente, usando la mappa definita sopra per passare da $(i,j)$ a $m$ attraverso $n$ (poiché il file txt contiene la coppia $m,n$).

Complessità: $O(N^2)$, come sopra.

\item \texttt{GeneraEsportaMatriceRhs}.
Genera i coefficienti non nulli della matrice $A$ del sistema ($N^2 \times N^2)$, il termine noto, e li esporta.
Prende in input:
\begin{enumerate}
\item $N$.
\item Una stringa "nat" o "ord" che decide il tipo di ordinamento (per dare il nome giusto ai file).
\item La mappe (matrice e vettore) generate dalle funzioni sopra (devono essere quelle relative all'ordinamento scelto).
\item La funzione che fa da sorgente di calore, che sarà definita nel main come funzione lambda.
Per darla come parametro, serve includere $\texttt{functional}$.
\item La diffusività $\kappa$, che di default è $0.01$.
\item Un vettore di $4N+4$ componenti che contiene le condizioni al bordo $u_0(x,y)$.
Di default la condizione di bordo dev'essere 0, ma non posso mettere come valore di default un vettore la cui lunghezza dipenda da $N$, che è un argomento della funzione.
Perciò di default il vettore dei valori al bordo è vuoto, e in tal caso lo ridefinisco all'interno della funzione in modo che contenga solo zeri.
\end{enumerate}

La funzione crea un file $\texttt{A.txt}$ e un file $\texttt{rhs.txt}$ (ne servirà uno per ogni $N$ e per i due ordinamenti).

Scorre sugli indici di riga $l$ di $A$, cioè sulle $N^2$ righe del sistema lineare.
Usando una delle mappe definite sopra, associa a $l$ una coppia $(i,j)$.
Calcola il valore $f(x_i,y_j)$ e lo salva in una variabile $b$, che sarà il coefficiente $l$ del termine noto.
Dall'equazione scritta nel testo si vede che solo 5 coefficienti ($(i+1,j)$ etc) sono coinvolti nella riga $(i,j)$ del sistema.
Usando una delle mappe definite sopra, associa a quelli che non sono di bordo (cioè hanno immagine non negativa tramite la mappa; sono almeno 3 e al massimo 5) un indice $m$.
Calcola il valore di $A_{l,m}$ e scrive sul file $\texttt{A.txt}$ la riga $l,m,A_{l,m}$.
Il formato è COO, cioè sto salvando indici e valore dei coefficienti non nulli.

Per gli eventuali indici di bordo, usa l'indice $m$ calcolato con la mappa per prendere il valore $u_0(x,y)$ assegnato, e lo somma alla variabile $b$ (con opportuno segno e coefficiente).
Infine scrive sul file $\texttt{rhs.txt}$ il valore di $b$.

Complessità: $O(N^2)$, perché il ciclo su $l$ fa $N^2$ iterazioni, e ogni iterazione contiene poche operazioni che costano $O(1)$ (il calcolo della funzione $f(x,y)$ potrebbe essere pesante, ma non dipende da $N$).

\end{itemize}

Complessità totale: $O(N^2)$.

\section{Fattorizzazione di Cholesky}
Input: matrice $A$ del sistema lineare.

Output: matrice $L$ tale che $A = -LL^T$.

File coinvolti: notebook $\texttt{soluzione\_sistema.ipynb}$ che legge $\texttt{A.txt}$ generato nel modulo precedente.

Strutture dati:
\begin{itemize}
\item array di NumPy, per salvare $A$ in formato COO. 
\item Formato CSC fornito da SciPy.
\end{itemize}

Funzioni:
\begin{itemize}
\item $\texttt{ImportaMatriceCSC}$.
Prende in input $N$ e una stringa "ord" o "nat" per scegliere il file.
Legge $\texttt{A.txt}$ in formato COO.
Con $\texttt{np.loadtxt}$ scarica le tre colonne del file come tre array che contengono indici di riga, indici di colonna e coefficienti, e converte in CSC (con $\texttt{scipy.sparse.csc\_matrix}$).
Restituisce la matrice $A$ in CSC.

Complessità: $O(N^2)$, che è il numero di entrate non nulle in $A$. Anche la conversione in CSC eseguita da SciPy dovrebbe richiedere $O(N^2)$.

\item $\texttt{CalcolaCholesky}$.
Prende in input una matrice $A$ (definita positiva; nel nostro caso sarà $-A$) e calcola $L$ tale che $LL^T = A$ con la funzione fornita nel testo.
Non sono sicuro della complessità della funzione usata.

\end{itemize}

\section{Soluzione del sistema lineare}
Input: fattorizzazione di $A$ e termine noto del sistema lineare.

Output: soluzione del sistema lineare.
Analisi del tempo necessario al variare di $N$ e dell'ordinamento dei nodi.
Plot della soluzione.

File coinvolti: notebook $\texttt{soluzione\_sistema.ipynb}$ che legge il file $\texttt{rhs.txt}$ generato nel modulo "approssimazione".

Strutture dati:
\begin{itemize}
\item array di NumPy per i vettori del termine noto e della soluzione, e per le matrici necessarie a passare da $n$ o $m$ a $(i,j)$.
\end{itemize}

Funzioni:
\begin{itemize}
\item $\texttt{ImportaRhs}$.
Legge $\texttt{rhs.txt}$ e scarica il termine noto come array $b$.
Complessità: $O(N^2)$, che è il numero di coefficienti di $b$.

\item $\texttt{RisolviSistema}$.
Prende in input $L$ triangolare inferiore, e $b$ array.
Risolve prima il sistema triangolare inferiore $Ly = -b$, poi quello triangolare superiore $L^T u= y$, con la funzione $\texttt{scipy.sparse.linalg.spsolve\_triangular}$.
L'output è il vettore $u$, nell'ordine scelto.
Complessità: $O(nnz(L))$, perché i due sistemi sono triangolari. $nnz(L) = O(N^3)$ per l'ordinamento naturale (in questo caso $A$ è a banda con larghezza $N$, la fattorizzazione dovrebbe conservare la banda ma potenzialmente la riempie tutta), $O(N^2 \log N)$ per il riordinamento.


\item $\texttt{GeneraMappaNaturale}$.
Legge il file $\texttt{coords.txt}$ e genera un array che in posizione $n$ (indice progressivo naturale) contiene la coppia $(i,j)$ associata. Complessità: $O(N^2)$, la lunghezza dell'array.

\item $\texttt{GeneraMappaOrdinata}$.
Legge il file $\texttt{ordering.txt}$ e genera un array che in posizione $m$ (indice progressivo dopo il riordinamento) contiene la coppia $(i,j)$ associata.
Per convertire da $m$ a $(i,j)$, passo attraverso $n$, usando la funzione precedente.
Uso il modo nativo di NumPy per permutare gli elementi ("fancy indexing"): valuto l'array "mappa naturale" sull'intero array di conversione tra gli ordinamenti $m$ e $n$.
Il risultato è l'array voluto che mappa $m$ in $(i,j)$.
Complessità: $O(N^2)$, perché scaricare dal file, chiamare la funzione precedente e permutare i coefficienti richiedono tutti $O(N^2)$.

Queste due funzioni sono analoghe alle funzioni $\texttt{GeneraMappaNaturale/Ordinata}$ del modulo "approssimazione lineare", ma operano solo in un verso (da $n$ o $m$ a $(i,j)$).

\item \texttt{GeneraMappaBordo}.
Genera un array di $4N+4$ elementi che in posizione $k$ contiene il $k+1$-esimo punto $(i,j)$ che si trova scorrendo in senso orario sul bordo del quadrato, partendo dall'origine.
Per farlo uso gli array di NumPy, che sono più veloci dei cicli for di Python. Complessità: $O(N)$, che è la lunghezza del bordo.

Questa funzione fa l'operazione inversa della funzione $\texttt{GeneraIndiciBordo}$ del modulo "approssimazione lineare", a meno del segno che lì era stato usato per distinguere tra punti di bordo e non.

\item $\texttt{DisegnaSoluzione}$.
Prende in input $N$, una stringa che specifica l'ordinamento, il vettore soluzione $u$ in quell'ordinamento e il vettore delle condizioni al bordo nel formato definito nel modulo "approssimazione lineare".
Usa le mappe definite sopra per trasformare dall'ordinamento scelto agli indici $(i,j)$, e disegna la soluzione.

\end{itemize}

\section{Analisi sperimentale}
\subsection{Struttura delle matrici $A,L$}
Nel notebook sono riportate le matrici $A,L$ nel caso $N=128$.

Con l'ordinamento naturale, $A$ è a banda simmetrica; sembra una banda completamente piena, ma in realtà in ogni riga ci sono al massimo 5 coefficienti non nulli.
Il fattore $L$ corrispondente è a sua volta a banda (inferiore), ma non è sparso, come si vedrà sotto.

Dopo il riordinamento, $A$ ha una struttura ricorsiva, con blocchi vuoti tra gli indici di $V_1$ e $V_2$ ai vari livelli della ricorsione.
Il fattore $L$ ha un numero di coefficienti non nulli molto minore rispetto all'ordinamento naturale.

\subsection{Grafico della soluzione}
Nel notebook è riportato il grafico della soluzione con $N=256$.
I bordi sono costretti a $T=0$, ma come atteso la temperatura sale vicino alla sorgente di calore localizzata intorno all'origine.

\subsection{Analisi dei tempi}
Con le funzioni $\texttt{ImportaMatrice}$, $\texttt{ImportaRhs}$, $\texttt{CalcolaCholesky}$, $\texttt{RisolviSistema}$ del modulo "soluzione", risolvo il sistema lineare al variare di $N$ tra 32 e 1024, e al variare dell'ordinamento (naturale e riordinato), con un ciclo for.
Il caso di $N=1024$ con ordinamento naturale ha richiesto troppo tempo e non è stato completato.

Misuro il tempo necessario a svolgere la fattorizzazione di Cholesky.
In un grafico bilogaritmico di $t$ in funzione di $N$, si osserva una pendenza di circa 3.2 nel caso naturale, 2.2 (escludendo il primo punto un po' anomalo) nel caso riordinato.
Non coincide del tutto con le attese $O(N^4)$ e $O(N^3)$.
Forse non siamo ancora nel limite asintotico, oppure la funzione usata riesce a ottimizzare il calcolo.

Misuro poi il tempo necessario per risolvere il sistema lineare.
Nel caso naturale la pendenza è circa 3 come atteso.
Nel caso riordinato non dovrebbe essere una legge di potenza ma comunque la pendenza è approssimativamente 2.4, intermedia tra 2 e 3 come atteso.

Il tempo necessario a risolvere il sistema dovrebbe essere circa proporzionale al numero di coefficienti non nulli in $L$, e in effetti dal grafico di $nnz(L)$ in funzione di $N$ si trovano circa le stesse pendenze. 

\end{document}
